\documentclass[12pt,a4paper]{article}

% ================== GÓI LỆNH ==================
\usepackage{tikz}
\usepackage[T5]{fontenc}
\usepackage[utf8]{inputenc}
\usepackage[vietnamese]{babel}
\usepackage{amsmath,amssymb}
\usepackage{geometry}
\usepackage{xcolor}
\usepackage{tcolorbox}
\usepackage{fancyhdr}
\usepackage{tocloft}
\usepackage{ragged2e}
\usepackage{setspace}
\usepackage[hidelinks]{hyperref}
\usepackage{bookmark}
\usepackage{lastpage}
\tcbuselibrary{skins,breakable}
\geometry{a4paper, margin=2.5cm}
\setstretch{1.2}
\usepackage{graphicx}
\usepackage[useregional]{datetime2}
% ================== THIẾT KẾ ĐẲNG CẤP - FONT ĐỒNG NHẤT ==================
\usepackage{fontawesome5}

\pagestyle{fancy}
\fancyhf{}
\renewcommand{\headrulewidth}{0pt}
\renewcommand{\footrulewidth}{0pt}

% --- HEADER: Tối giản, không nền ---
\fancyhead[L]{
    \begin{tikzpicture}[remember picture, overlay]

        % Đường kẻ ngang màu tím nhấn dưới Header (chừa lề 1cm)
        \draw[mainpurple, line width=0.8pt] 
        ([yshift=-2.2cm, xshift=1cm]current page.north west) -- 
        ([yshift=-2.2cm, xshift=-1cm]current page.north east);
        
        % Nội dung Header L
        \node[anchor=west, font=\small\color{darkpurple}] 
        at ([yshift=-1.65cm, xshift=1cm]current page.north west) {
            \faLayerGroup \hspace{6pt} \textbf{Hình học cao cấp hai chiều và ba chiều}
        };
    \end{tikzpicture}
}

\fancyhead[R]{
    \begin{tikzpicture}[remember picture, overlay]
        % Nội dung Header R
        \node[anchor=east, font=\small\color{darkpurple}] 
        at ([yshift=-1.65cm, xshift=-1cm]current page.north east) {
            \faCalendarAlt \hspace{4pt} \textbf{Năm học: 2025 - 2026}
        };
    \end{tikzpicture}
}

% --- FOOTER: Tối giản, chuyên nghiệp (Font đồng bộ thân trang) ---
\fancyfoot[L]{
    \begin{tikzpicture}[remember picture, overlay]
        % Đường kẻ mờ nền chạy ngang tinh tế
        \draw[mainpurple!30, line width=0.6pt] ([yshift=1.5cm, xshift=1cm]current page.south west) -- ([yshift=1.5cm, xshift=-1cm]current page.south east);
        
        % Thông tin sinh viên
        \node[anchor=west, font=\footnotesize\color{darkpurple}] at ([yshift=1.1cm, xshift=1cm]current page.south west) {
            \faIdBadge \hspace{4pt} Cẩn thận với người im lặng vì họ thường không nói gì
        };
    \end{tikzpicture}
}

\fancyfoot[R]{
    \begin{tikzpicture}[remember picture, overlay]
        % Số trang
        \node[anchor=east, font=\small\color{darkpurple}] at ([yshift=1.1cm, xshift=-1cm]current page.south east) {
            \faChevronRight \hspace{3pt} Trang \textbf{\thepage} \scalebox{0.8}
        };
    \end{tikzpicture}
}

% Điều chỉnh lề tối ưu
\geometry{
    a4paper,
    left=2.5cm,
    right=2.5cm,
    top=3cm,
    bottom=2.8cm,
    headheight=16pt,
    headsep=18pt,
    footskip=25pt
}
% ===========================================================================

% ================== MÀU SẮC ==================
\definecolor{mainpurple}{HTML}{5D32A8} % Màu tím đậm
\definecolor{bgpurple}{HTML}{F8F7FC}   % Màu nền nhạt

\definecolor{mainpurple}{RGB}{91, 45, 145}
\definecolor{bgpurple}{RGB}{245, 240, 255}
\definecolor{darkpurple}{RGB}{60, 30, 100}

\newtcolorbox{BaiToanBox}[1]{
    enhanced,
    breakable,
    sharp corners,
    boxrule=0pt,
    frame hidden,
    colback=bgpurple,
   borderline west={6pt}{0pt}{mainpurple!50}, % Viền trái đậm và dày hơn {mainpurple!50}, % Thanh dọc nhạt hơn chút
    top=10pt,
    bottom=10pt,
    left=15pt,
    right=10pt,
    overlay={
        % 1. Miếng gấp nhỏ phía dưới nhãn (tạo hiệu ứng 3D bên trái)
        \fill[darkpurple] ([xshift=-4pt, yshift=0pt]frame.north west) -- ++(4pt,0) -- ++(0,-6pt) -- cycle;
        
        % 2. Thân nhãn chính (nhô ra trái 4pt)
        \fill[mainpurple] 
            ([xshift=-4pt, yshift=0pt]frame.north west) -- 
            ++(0, 20pt) -- 
            ++(95pt, 0pt) -- 
            ++(14pt, -10pt) -- % Mũi nhọn đuôi
            ++(-14pt, -10pt) -- 
            cycle;
        
        % 3. Chữ "Bài toán"
        \node[anchor=west, text=white, font=\bfseries\large] 
            at ([xshift=6pt, yshift=10pt]frame.north west) {Bài tập #1};
            
        % 4. Hai dấu mũi tên >> (căn chỉnh sát đuôi vát)
        \begin{scope}[shift={([xshift=103pt, yshift=10pt]frame.north west)}, scale=1]
            \draw[mainpurple, line width=2.8pt, line cap=round] 
                (0.1, 0.35) -- (0.45, 0) -- (0.1, -0.35);
            \draw[mainpurple, line width=2.8pt, line cap=round] 
                (0.4, 0.35) -- (0.75, 0) -- (0.4, -0.35);
        \end{scope}
    }
}
% Lệnh sử dụng nhanh
\newcommand{\Cau}[2]{
    \vspace{0.8cm} % Tạo khoảng trống phía trên để nhãn không bị dính vào đoạn trên
    \begin{BaiToanBox}{#1}
        #2
    \end{BaiToanBox}
}
\newtcolorbox{DinhNghiaBox}[2][]{
    enhanced,
    before skip=10pt,after skip=10pt,
    colback=white, % Màu nền hộp
    colframe=mainpurple, % Màu viền (đang dùng màu tím của bạn, có thể đổi sang xanh blue nếu muốn giống ảnh 100%)
    boxrule=1pt,
    sharp corners,
    fonttitle=\bfseries,
    colbacktitle=white,
    attach boxed title to top left={xshift=0mm, yshift=0mm},
    boxed title style={
        sharp corners,
        size=small,
        colback=mainpurple, % Màu nền cái nhãn
        colframe=mainpurple,
        top=2pt, bottom=2pt, left=5pt, right=15pt, % Tạo khoảng trống bên phải để vát
    },
    title={#2},
    % Vẽ cái đuôi vát chéo bằng code TikZ
    overlay unbroken and first={
        \fill[mainpurple] (title.north east) -- ([xshift=10pt]title.south east) -- (title.south east) -- cycle;
    },
    #1
}
% ================== MỤC LỤC CHUYÊN NGHIỆP ==================

\renewcommand{\contentsname}{\textcolor{mainpurple}{\Huge\bfseries MỤC LỤC}}

% Canh giữa tiêu đề
\renewcommand{\cfttoctitlefont}{\hfill\Huge\bfseries\color{mainpurple}}
\renewcommand{\cftaftertoctitle}{\hfill}

% Khoảng cách sau tiêu đề
\setlength{\cftaftertoctitleskip}{15pt}

% ===== Section =====
\renewcommand{\cftsecfont}{\bfseries}
\renewcommand{\cftsecpagefont}{\bfseries}

\setlength{\cftbeforesecskip}{8pt}

% ===== Subsection =====
\setlength{\cftbeforesubsecskip}{2pt}

% ===== Dot leader đẹp =====
\renewcommand{\cftdotsep}{1.5}

% ===== Số mục đậm hơn =====
\renewcommand{\cftsecleader}{\cftdotfill{\cftdotsep}}

% ===== Thụt lề mục =====
\setlength{\cftsecindent}{0pt}
\setlength{\cftsubsecindent}{1.5em}

% ===== Độ rộng số mục =====
\setlength{\cftsecnumwidth}{2.5em}


% =====================================================
\begin{document}

\pagestyle{empty}

\begin{tikzpicture}[remember picture,overlay]

% Khung ngoài
\draw[line width=2pt, color=mainpurple]
([xshift=2cm,yshift=-2cm]current page.north west)
rectangle
([xshift=-2cm,yshift=2cm]current page.south east);

% Khung trong
\draw[line width=0.8pt, color=mainpurple!70]
([xshift=2.3cm,yshift=-2.3cm]current page.north west)
rectangle
([xshift=-2.3cm,yshift=2.3cm]current page.south east);

\end{tikzpicture}

\vspace*{0.3cm}

\begin{center}
\begin{minipage}{0.9\textwidth}

\vspace{10pt}

\begin{center}
{\Large\textbf{TRƯỜNG ĐẠI HỌC SƯ PHẠM THÀNH PHỐ HỒ CHÍ MINH}}\\[6pt]
{\large\textbf{KHOA TOÁN - TIN HỌC}}
\end{center}

\vspace{0.4cm}

\begin{center}
\includegraphics[width=0.45\textwidth]{logohcmue.png}\hspace{1cm}
\includegraphics[width=0.25\textwidth]{Logo KhoaToan.png}
\end{center}

\vspace{0.35cm}

\begin{center}
\rule{0.9\textwidth}{1pt}\\[6pt]
{\LARGE\bfseries HÌNH HỌC CAO CẤP \\ HAI CHIỀU VÀ BA CHIỀU}\\[6pt]
\rule{0.9\textwidth}{1pt} \\[15pt]
\Large\textbf{BÀI TẬP NHÓM}
\end{center}

\vspace{0.2cm}

\begin{center}

\textbf{NHÓM 6 - Lớp: 51.01.TUD}

\vspace{0.2cm}

\begin{tabular}{l}
• Nguyễn Phạm Lê Vy - 51.01.108.052 \\
• Trương Bảo Chí - 51.01.108.008 \\
• Dương Nguyên Khoa - 51.01.108.023 \\
• Nguyễn Thị Thanh Thảo - 51.01.108.039\\
• Nguyễn Huỳnh Trúc Thy - 51.01.108.044\\
• Nguyễn Ngọc Thuỷ Tiên - 51.01.108.046
\end{tabular} \\[12pt]

\textbf{GIẢNG VIÊN HƯỚNG DẪN: ThS. PHAN NGỌC YẾN}\\[6pt]
Lớp học phần: 2521APMA180701

\end{center}

\vspace{10pt}

\end{minipage}
\end{center}


\newpage
\pagestyle{fancy}

% ================== MỤC LỤC ==================
\tableofcontents
\newpage

% =====================================================
\section{Bài 1. Không gian Affine và Hình học Affine. \\[4pt] Bài 2. Tọa độ Affine và mục tiêu Affine}

\Cau{2}{ Cho hai không gian affine $(\mathbb{A}_{\nVdash}, \varphi_1, V_1)$,
$(\mathbb{A}_{\nvDash}, \varphi_2, V_2)$ trên cùng một trường $\mathbb{K}$.
Đặt
\(
\mathbb{A} = \mathbb{A}_{\nVdash} \times \mathbb{A}_{\nvDash}, \text{ }
V = V_1 \times V_2
\)
và \( \varphi: \mathbb{A} \times \mathbb{A} \to V \) là ánh xạ cho bởi:
\[
\varphi\big((M_1, M_2), (N_1, N_2)\big)
:= \big(\varphi_1(M_1, N_1), \varphi_2(M_2, N_2)\big).
\]
Chứng minh rằng bộ ba $(\mathbb{A},\varphi,V)$ là một không gian Affine.
}
\begin{center}
{\large\bfseries\textcolor{mainpurple}{Lời giải}}
\end{center}
\begin{DinhNghiaBox}{Ghi chú}
Bộ ba \( (\mathbb{A}, \varphi, \mathbb{V}) \) được gọi là không gian Affine nếu hai tiên đề sau đây thỏa mãn:
\begin{itemize}
    \item Với mọi điểm \( M \) thuộc \( \mathbb{A}\) và mọi vectơ \( \vec{u} \) thuộc \( \mathbb{V} \) có duy nhất điểm \(N\) thuộc \( \mathbb{A} \) sao cho \( \overrightarrow{MN} = \vec{u} \)
    \item Với mọi ba điểm \(M,N,P\) thuộc \( \mathbb{A} \) ta luôn có: \( \overrightarrow{MN} + \overrightarrow{NP} = \overrightarrow{MP} \)
\end{itemize}
\end{DinhNghiaBox}
\textbf{Kiểm tra tiên đề (i):} \\[4pt]
Giả sử \( M=(M_1,M_2) \in \mathbb{A} \) và \( \vec{v}=(\vec{v_1},\vec{v_2}) \in V\) \\[4pt]
Do $(\mathbb{A}_{\nVdash}, \varphi_1, V_1) $ là không gian affine nên với \(M_1 \in \mathbb{A}_{\nVdash}, \vec{v_1} \in V_1 \), tồn tại duy nhất một điểm \( N_1 \in \mathbb{A}_{\nVdash} \) thỏa mãn \( \varphi_1(M_1,N_1)=\vec{v_1} \) \\[4pt]
Do $(\mathbb{A}_{\nvDash}, \varphi_2, V_2) $ là không gian affine nên với \(M_2 \in \mathbb{A}_{\nvDash}, \vec{v_2} \in V_2 \), tồn tại duy nhất một điểm \( N_2 \in \mathbb{A}_{\nvDash} \) thỏa mãn \( \varphi_2(M_2,N_2)=\vec{v_2} \) \\[4pt]
Ta có: \[
\varphi \left ( (M_1,M_2),(N_1,N_2) \right ) = \left ( \varphi_1 (M_1,N_1),\varphi_2 (M_2,N_2) \right ) = (\vec{v_1},\vec{v_2})=\vec{v}
\]
Do đó, tồn tại duy nhất một điểm \(N=(N_1,N_2) \in \mathbb{A}\) sao cho \( \varphi(M,N)=\vec{v} \) nên tiên đề (i) thỏa mãn \qquad (1) \\[4pt]
\textbf{Kiểm tra tiên đề (ii):} \\[4pt]
Giả sử $M = (M_1, M_2)$, $N = (N_1, N_2)$ và $P = (P_1, P_2)$ thuộc \( \mathbb{A}\) \\[4pt]
Ta có: \[
\begin{aligned}
\varphi(M,N) + \varphi(N,P)
&= \varphi\bigl((M_1,M_2),(N_1,N_2)\bigr) + \varphi\bigl((N_1,N_2),(P_1,P_2)\bigr)  \\
&= \bigl(\varphi_1(M_1,N_1),\varphi_2(M_2,N_2)\bigr) + \bigl(\varphi_1(N_1,P_1),\varphi_2(N_2,P_2)\bigr)  \\
&= \bigl(\varphi_1(M_1, N_1) + \varphi_1(N_1, P_1),\, \varphi_2(M_2, N_2) + \varphi_2(N_2, P_2)\bigr). \\
&= (\varphi_1(M_1, P_1), \varphi_2(M_2, P_2)) \\
&= \varphi((M_1, M_2), (P_1, P_2)) \\
&= \varphi(M, P)
\end{aligned}
\]
Suy ra tiên đề (ii) thỏa mãn \qquad (2) \\[4pt]
Từ (1) và (2) suy ra: Bộ ba $(\mathbb{A},\varphi,V)$ là một không gian Affine.

\section{Bài 3. Các phẳng trong không gian Affine}
\Cau{6}{ Trong không gian affine $\mathbb{A}^3$, cho mục tiêu affine 
$\{O; \vec e_1, \vec e_2, \vec e_3\}$ và các điểm $P_i$ với 
$\overrightarrow{OP_i} = a_i \vec e_i$, $a_i \ne 0$, $i = 1,2,3$. Chứng minh rằng ba điểm $P_1, P_2, P_3$ độc lập và phương trình siêu phẳng 
đi qua ba điểm ấy có thể viết dưới dạng:
\[
\frac{x_1}{a_1} + \frac{x_2}{a_2} + \frac{x_3}{a_3} = 1.
\] }
\begin{center}
{\large\bfseries\textcolor{mainpurple}{Lời giải}}
\end{center}
\setlength{\fboxsep}{6pt}
\fbox{\textbf{\large Cách 1: Dựa trên giáo trình Hình học tuyến tính}}
\begin{DinhNghiaBox}{Ghi chú}
Hệ \(m+1\) điểm \(P_0,P_1,...,P_m\) \((m \ge 1)\) của không gian Affine \(\mathbb{A}\) được gọi là độc lập nếu \(m\) vectơ \( \overrightarrow{P_0P_1},\overrightarrow{P_0P_2},...,\overrightarrow{P_0P_m} \) của không gian vectơ \( \vec{\mathbb{A}} \) là hệ vectơ độc lập tuyến tính 
\end{DinhNghiaBox} 
\noindent Xét \( \overrightarrow{P_1P_j}=\overrightarrow{OP_j}-\overrightarrow{OP_1}=a_j\vec{e_j}-a_1\vec{e_1} \) với \( j= 2,3 \) \\[4pt]
Ta có: \[
\sum_{j=2}^{3} x_j \overrightarrow{P_1P_j} = \vec{0} 
\Leftrightarrow \sum_{j=2}^{3} x_j \left ( a_j\vec{e_j}-a_1\vec{e_1} \right ) = \vec{0}
\Leftrightarrow \sum_{j=2}^{3} x_j a_j\vec{e_j}+\sum_{j=2}^{3} (-x_j a_1)\vec{e_1}=\vec{0}
\]
Do $\{\vec e_1, \vec e_2, \vec e_3\}$ độc lập tuyến tính nên \(
\begin{cases}
x_j a_j = 0,  \forall j=2,3 \\[6pt]
\displaystyle \sum_{j=2}^{3} (-x_j a_1) = 0.
\end{cases}
\) \\[4pt]
Mà \( a_i \neq 0, \forall i=1,2,3 \) nên \( x_j=0, \forall j=2,3 \). Suy ra: \( \overrightarrow{P_1P_2},\overrightarrow{P_1P_3} \) độc lập tuyến tính (độc lập Affine). Do đó ba điểm \(P_1,P_2,P_3\) độc lập
\begin{DinhNghiaBox}{Ghi chú}
\begin{itemize}
    \item Qua \(m+1\) điểm của không gian affine \(\mathbb{A}\) có một và chỉ một \(m\)-phẳng \((m \ge 0)\)
    \item Nếu \( \dim \mathbb{A}=n\) thì \((n-1)\)-phẳng được gọi là siêu phẳng của không gian đó
\end{itemize}
\end{DinhNghiaBox}
\noindent Suy ra tồn tại duy nhất 2-phẳng qua \(P_1,P_2,P_3\) \\[4pt]
Suy ra tồn tại duy nhất siêu phẳng qua \(P_1,P_2,P_3\) \\[4pt]
Xét siêu phẳng \( \dfrac{x_1}{a_1} + \dfrac{x_2}{a_2} + \dfrac{x_3}{a_3} = 1 \quad (\alpha) \) \\[4pt]
Thay tọa độ ba điểm \( P_1, P_2, P_3 \) vào phương trình siêu phẳng \( (\alpha) \) , ta thấy ba điểm đó thuộc siêu phẳng đó. Vì vậy \( (\alpha) \) là siêu phẳng duy nhất chứa ba điểm \( P_1, P_2, P_3 \). \\[4pt]
Vậy: Phương trình siêu phẳng đi qua ba điểm \(P_1,P_2,P_3\) có thể viết dưới dạng \[ \dfrac{x_1}{a_1} + \dfrac{x_2}{a_2} + \dfrac{x_3}{a_3} = 1 \]
\setlength{\fboxsep}{6pt}
\fbox{\textbf{\large Cách 2: Trình bày theo gợi ý của A.I}} \\[4pt]
Đối với mục tiêu affine \( \{O; \vec e_1, \vec e_2, \vec e_3\}  \), giả sử \( P_1=(a_1;0;0),P_2=(0;a_2;0),P_3=(0;0;a_3) \) \\[4pt]
Ta có: \( \overrightarrow{P_1P_2} =(-a_1;a_2;0)\) và \( \overrightarrow{P_1P_3} =(-a_1;0;a_3)\) \\[4pt]
Xét ma trận các hệ số của hệ trên là: \(
A=
\begin{bmatrix}
 -a_1 & a_2 & 0\\
 -a_1 & 0 & a_3
\end{bmatrix}
\;\xrightarrow{\,d_2 \rightarrow -d_1+d_2\,}\;
\begin{bmatrix}
 -a_1 & a_2 & 0\\
 0 & -a_2 & a_3
\end{bmatrix}
 \) \\[4pt]
Ta thấy: \( rankA=2= \) số vectơ nên hệ vectơ \( \overrightarrow{P_1P_2} \), \( \overrightarrow{P_1P_3} \) độc lập tuyến tính \\[4pt]
Suy ra: Ba điểm \( P_1, P_2, P_3 \) độc lập \\[4pt]
Siêu phẳng đi qua ba điểm $P_1, P_2, P_3$ trong $\mathbb{A}^3$ có phương trình tổng quát dạng: $$Ax_1 + Bx_2 + Cx_3 + D = 0 \qquad (*)$$
Thay tọa độ điểm \( P_1=(a_1;0;0) \) vào phương trình (*), ta có: \(
A.a_1+D=0 \Leftrightarrow A=\dfrac{-D}{a_1}
\) \\[4pt]
Thay tọa độ điểm \( P_2=(0;a_2;0) \) vào phương trình (*), ta có: \(
B.a_2+D=0 \Leftrightarrow B=\dfrac{-D}{a_2}
\) \\[4pt]
Thay tọa độ điểm \( P_3=(0;0;a_3) \) vào phương trình (*), ta có: \(
C.a_3+D=0 \Leftrightarrow C=\dfrac{-D}{a_3}
\) \\[4pt]
Khi đó, phương trình (*) tương đương: \[
\dfrac{-D}{a_1}x_1+\dfrac{-D}{a_2}x_2+\dfrac{-D}{a_3}x_3+D=0 \quad (1)
\]
Do ba điểm \( P_1,P_2,P_3 \) độc lập nên phương trình (*) không đi qua gốc tọa độ \( O\), suy ra \(D \neq 0 \) \\[4pt]
Chia hai vế cho D, phương trình (1) trở thành: \[
\dfrac{-1}{a_1}x_1+\dfrac{-1}{a_2}x_2+\dfrac{-1}{a_3}x_3+1=0 \Leftrightarrow \frac{x_1}{a_1} + \frac{x_2}{a_2} + \frac{x_3}{a_3} = 1 \quad \text{(đpcm)}
\]

\section{Bài 4. Tâm tỉ cự của một hệ điểm \\[4pt] Bài 5. Tập lồi trong không gian Affine thực \\[4pt] Bài 6. Ánh xạ affine của các không gian affine và phép biến đổi của không gian affine}
\Cau{10}{ Trong không gian affine \(\mathbb{A}^n\) cho công thức đổi mục tiêu: 
\[
\begin{cases}
x_1 = y_1 + y_2 + y_3 + \cdots + y_n + 1 \\
x_2 = \qquad y_2 + y_3 + \cdots + y_n + 2 \\
x_3 = \qquad\qquad y_3 + \cdots + y_n + 3 \\
\cdots \\
x_n = \qquad\qquad\qquad\qquad y_n + n
\end{cases}
\]
trong đó \((x_1,x_2,...,x_n)\) tọa độ của một điểm \(X\) thuộc \(\mathbb{A}^n\) đối với mục tiêu thứ nhất và \((y_1,y_2,...,y_n)\) là tọa độ của điểm \(X\) đối với mục tiêu thứ hai. Lập phương trình phép biến đổi affine đối với mục tiêu thứ nhất và biến mục tiêu thứ nhất thành mục tiêu thứ hai. Tìm tọa độ các đỉnh của mục tiêu thứ hai đối với mục tiêu thứ nhất. }
\begin{center}
{\large\bfseries\textcolor{mainpurple}{Lời giải}}
\end{center}
\begin{DinhNghiaBox}{Ghi chú}
Ta đã biết công thức đổi mục tiêu \( [x]=C^*[x']+[a_0] \), trong đó \([x], [x'] \) theo thứ tự là tọa độ của điểm X đối với mục tiêu \( \left \{ E_0; E_i \right \} \) và đối với mục tiêu \( \left \{ E'_0; E'_i \right \} \). Còn C là ma trận chuyển từ mục tiêu \( \left \{ E_0; E_i \right \} \) sang \( \left \{ E'_0; E'_i \right \} \).
\end{DinhNghiaBox}
\textbf{a) Lập phương trình phép biến đổi affine đối với mục tiêu thứ nhất và biến mục tiêu thứ nhất thành mục tiêu thứ hai} \\[4pt]
Theo giả thiết của đề bài, ta suy ra phương trình phép biến đổi affine đối với mục tiêu thứ nhất \( \left \{ E_0; E_i \right \} \) và biến mục tiêu thứ nhất thành mục tiêu thứ hai \( \left \{ E'_0; E'_i \right \} \) là: \[
[x']=C^*[x]+[a_0]
\]
\begin{center}
    \begin{tikzpicture}[scale=0.9]

% Hai hệ cơ sở
\node (E)  at (0,0) {$\{E_0,E_i\}$};
\node (Ep) at (6,0) {$\{E'_0,E'_i\}$};

% Điểm x
\coordinate (X) at (6,4);
\node[above] at (X) {$X$};

% Mũi tên ánh xạ
\draw[->,thick] (E) -- (Ep)
node[midway,above] {$f$}
node[midway,below] {$C$};

% Đường nét đứt (không đụng chữ)
\draw[dashed,thick] ([yshift=6pt]E) -- (X);
\node at (2.2,2.2) {$[x]$};

\draw[dashed,thick] ([yshift=6pt]Ep) -- (X);
\node[right] at (6,2) {$[x']$};

\end{tikzpicture}
\end{center}
Ta có:
\[
C^* =
\begin{bmatrix}
1 & 1 & 1 & \dots & 1 & 1 \\
0 & 1 & 1 & \dots & 1 & 1 \\
0 & 0 & 1 & \dots & 1 & 1 \\
\vdots & \vdots & \vdots & \ddots & \vdots \\
0 & 0 & 0 & \dots & 1 & 1 \\
0 & 0 & 0 & \dots & 0 & 1
\end{bmatrix},
\quad
a_0 =
\begin{bmatrix}
1 \\ 2 \\ 3 \\ \vdots \\ n-1 \\  n
\end{bmatrix}
\]
Phương trình phép biến đổi affine đối với mục tiêu thứ nhất \( \left \{ E_0; E_i \right \} \) và biến mục tiêu thứ nhất thành mục tiêu thứ hai \( \left \{ E'_0; E'_i \right \} \) là: \[
\begin{bmatrix}
    x'_1 \\ x'_2 \\ x'_3 \\ ... \\ x'_{n-1} \\ x'_n
\end{bmatrix} = \begin{bmatrix}
1 & 1 & 1 & \dots & 1 & 1 \\
0 & 1 & 1 & \dots & 1 & 1 \\
0 & 0 & 1 & \dots & 1 & 1 \\
\vdots & \vdots & \vdots & \ddots & \vdots \\
0 & 0 & 0 & \dots & 1 & 1 \\
0 & 0 & 0 & \dots & 0 & 1
\end{bmatrix} \begin{bmatrix}
    x_1 \\ x_2 \\ x_3 \\ ... \\ x_{n-1} \\ x_n
\end{bmatrix} + \begin{bmatrix}
    1 \\ 2 \\ 3 \\ ... \\ n-1 \\ n
\end{bmatrix}
\]
Hay \[ \begin{cases} x'_1 = x_1 + x_2 + x_3 + \dots + x_n + 1 \\ x'_2 = x_2 + x_3 + \dots + x_n + 2 \\ x'_3 = x_3 + \dots + x_n + 3 \\ \dots \\ x'_{n-1}=x_{n-1}+x_n+n-1 \\ x'_n = x_n + n \end{cases} \]
\textbf{b) Tìm tọa độ các đỉnh của mục tiêu thứ hai đối với mục tiêu thứ nhất } \\[4pt]
Đối với điểm \(E_0'\), ta có: 
\[
\begin{bmatrix}
    x'_1 \\ x'_2 \\ x'_3 \\ ... \\ x'_{n-1} \\ x'_n
\end{bmatrix} = \begin{bmatrix}
1 & 1 & 1 & \dots & 1 & 1 \\
0 & 1 & 1 & \dots & 1 & 1 \\
0 & 0 & 1 & \dots & 1 & 1 \\
\vdots & \vdots & \vdots & \ddots & \vdots \\
0 & 0 & 0 & \dots & 1 & 1 \\
0 & 0 & 0 & \dots & 0 & 1
\end{bmatrix} \begin{bmatrix}
    0 \\ 0 \\ 0 \\ \dots \\ 0 \\ 0
\end{bmatrix} + \begin{bmatrix}
    1 \\ 2 \\ 3 \\ ... \\ n-1 \\ n
\end{bmatrix} = \begin{bmatrix}
    1 \\ 2 \\ 3 \\ \dots \\ n-1 \\ n
\end{bmatrix}
\]
Đối với điểm \(E_1'\), ta có: 
\[
\begin{bmatrix}
    x'_1 \\ x'_2 \\ x'_3 \\ ... \\ x'_{n-1} \\ x'_n
\end{bmatrix} = \begin{bmatrix}
1 & 1 & 1 & \dots & 1 & 1 \\
0 & 1 & 1 & \dots & 1 & 1 \\
0 & 0 & 1 & \dots & 1 & 1 \\
\vdots & \vdots & \vdots & \ddots & \vdots \\
0 & 0 & 0 & \dots & 1 & 1 \\
0 & 0 & 0 & \dots & 0 & 1
\end{bmatrix} \begin{bmatrix}
    1 \\ 0 \\ 0 \\ \dots \\ 0 \\ 0
\end{bmatrix} + \begin{bmatrix}
    1 \\ 2 \\ 3 \\ ... \\ n-1 \\ n
\end{bmatrix} = \begin{bmatrix}
    2 \\ 2 \\ 3 \\ \dots \\ n-1 \\ n
\end{bmatrix}
\]
Đối với điểm \(E_2'\), ta có: 
\[
\begin{bmatrix}
    x'_1 \\ x'_2 \\ x'_3 \\ ... \\ x'_{n-1} \\ x'_n
\end{bmatrix} = \begin{bmatrix}
1 & 1 & 1 & \dots & 1 & 1 \\
0 & 1 & 1 & \dots & 1 & 1 \\
0 & 0 & 1 & \dots & 1 & 1 \\
\vdots & \vdots & \vdots & \ddots & \vdots \\
0 & 0 & 0 & \dots & 1 & 1 \\
0 & 0 & 0 & \dots & 0 & 1
\end{bmatrix} \begin{bmatrix}
    0 \\ 1 \\ 0 \\ \dots \\ 0 \\ 0
\end{bmatrix} + \begin{bmatrix}
    1 \\ 2 \\ 3 \\ ... \\ n-1 \\ n
\end{bmatrix} = \begin{bmatrix}
    2 \\ 3 \\ 3 \\ \dots \\ n-1 \\ n
\end{bmatrix}
\]
... \\[4pt]
Đối với điểm \(E_{n-1}'\), ta có: 
\[
\begin{bmatrix}
    x'_1 \\ x'_2 \\ x'_3 \\ ... \\ x'_{n-1} \\ x'_n
\end{bmatrix} = \begin{bmatrix}
1 & 1 & 1 & \dots & 1 & 1 \\
0 & 1 & 1 & \dots & 1 & 1 \\
0 & 0 & 1 & \dots & 1 & 1 \\
\vdots & \vdots & \vdots & \ddots & \vdots \\
0 & 0 & 0 & \dots & 1 & 1 \\
0 & 0 & 0 & \dots & 0 & 1
\end{bmatrix} \begin{bmatrix}
    0 \\ 0 \\ 0 \\ \dots \\ 1 \\ 0
\end{bmatrix} + \begin{bmatrix}
    1 \\ 2 \\ 3 \\ ... \\ n-1 \\ n
\end{bmatrix} = \begin{bmatrix}
    2 \\ 3 \\ 4 \\ \dots \\ n \\ n
\end{bmatrix}
\]
Đối với điểm \(E_n'\), ta có: 
\[
\begin{bmatrix}
    x'_1 \\ x'_2 \\ x'_3 \\ ... \\ x'_{n-1} \\ x'_n
\end{bmatrix} = \begin{bmatrix}
1 & 1 & 1 & \dots & 1 & 1 \\
0 & 1 & 1 & \dots & 1 & 1 \\
0 & 0 & 1 & \dots & 1 & 1 \\
\vdots & \vdots & \vdots & \ddots & \vdots \\
0 & 0 & 0 & \dots & 1 & 1 \\
0 & 0 & 0 & \dots & 0 & 1
\end{bmatrix} \begin{bmatrix}
    0 \\ 0 \\ 0 \\ \dots \\ 0 \\ 1
\end{bmatrix} + \begin{bmatrix}
    1 \\ 2 \\ 3 \\ ... \\ n-1 \\ n
\end{bmatrix} = \begin{bmatrix}
    2 \\ 3 \\ 4 \\ \dots \\ n \\ n+1
\end{bmatrix}
\]
Do đó tọa độ các đỉnh của mục tiêu thứ hai là: \[
\begin{aligned}
    E_0' &= (1,2,3,4,...,n-1,n) \\
    E_1' & = (2,2,3,4,...,n-1,n) \\
    E_2' & = (2,3,3,4,...,n-1,n) \\
    ... \\
    E'_{n-1} & = (2,3,4,5,...,n,n) \\
    E_n' & = (2,3,4,5,...,n,n+1)
\end{aligned}
\]

\end{document}
